% Provenance (never printed): see sources/outline_sources.md (Year 2: eccentricity, obliquity and
% precession; insolation computed from spherical astronomy; the Quaternary cycles; carbon dioxide
% and feedbacks). Orbital elements and tabulated insolation of Berger and Loutre (1991) via NOAA
% NCEI (data/derived/bachelor-2-orbit91, fetch_data Book 4 block); periods after Laskar et al.
% (2004) and the NASA Milankovitch page; history after the NASA Earth Observatory feature (2000);
% terminations after Barker and Knorr (2021); the mid-Pleistocene transition after Clark et al.
% (2006) and Koehler and van de Wal (2020). Figures from figdata/bachelor-2/orbital.py; values of
% the exercises and problem computed in the Book 4 scratch calculator calc/ch25.py.
\chapter{Orbital Forcing and the Ice Ages}\label{ch:b2:orbital-forcing}

Over the last million years the shape of the Earth's orbit has changed the sunlight received by the
whole planet in a year by less than two parts in a thousand: between \num{340.00} and
\qty{340.54}{W.m^{-2}} at the top of the atmosphere. Over the same million years the daily-mean
sunshine of mid-July at \ang{65} N has swung between about 388 and \qty{497}{W.m^{-2}}, by more than
a quarter. The ice ages followed the second number, not the first. This chapter computes both
from the elements of the Earth's orbit, and asks how sunlight at one latitude in one season can
pace the growth and decay of ice sheets that cover continents.
% ledger: data:b4-orbital, berger91:readme

\begin{recall}
\textbf{From the physics books}: the ellipse, its semi-major axis $a$ and eccentricity $e$, Kepler's
laws. \textbf{From the Year 1 volume}: the ecliptic and the obliquity $\varepsilon$, the Sun's
declination $\delta$ with $\sin\delta = \sin\varepsilon\sin\lambda$ ($\lambda$ the Sun's true
longitude, counted along the ecliptic from the vernal equinox), the tropical year, the daily-mean
insolation
\[
  \bar Q = \frac{S_0}{\pi}\left(\frac{a}{r}\right)^{2}\bigl(h_0\sin\varphi\sin\delta +
  \cos\varphi\cos\delta\sin h_0\bigr),\qquad \cos h_0 = -\tan\varphi\tan\delta ,
\]
and its global and annual mean, $S_0/(4\sqrt{1 - e^2})$. \textbf{From the school volumes}: the
glacial cycle and the Last Glacial Maximum; radiative forcing. \textbf{From earlier chapters}: the
\omterm{def:b2:precambrian-earth:ice-albedo}{ice--albedo feedback} (\cref{def:b2:precambrian-earth:ice-albedo}); the benthic stack, the EPICA
Dome C record and \omterm{def:b2:climate-proxies:age-model}{age models} (\cref{ch:b2:climate-proxies}).
\end{recall}
% ledger: earth:TSI, earth:obliquity, earth:ecc

\section{The orbit and its three elements}\label{sec:b2:orbital-forcing:elements}

The Earth's orbit is an ellipse with the Sun at one focus. The planets, Jupiter and Saturn above
all, slowly deform it and turn it in its plane; the pull of the Sun and the Moon on the equatorial
bulge makes the axis of rotation wobble like a top. Three quantities describe what matters for
sunlight: how elongated the ellipse is, how tilted the axis is, and where along the orbit the
seasons fall.
% ledger: nasa:milank-causes

\begin{definition}[Longitude of perihelion; climatic precession]\label{def:b2:orbital-forcing:elements}
The \emph{longitude of perihelion}\index{longitude of perihelion} $\varpi$ is the angle, measured in
the plane of the orbit from the direction of the vernal equinox, between that direction and the
direction of the perihelion seen from the Sun; it is counted from the moving equinox, so it
changes with both the wobble of the axis and the turning of the orbit. The \emph{climatic
precession}\index{climatic precession} is the quantity $e\sin\varpi$, also called the precession
index.
\end{definition}
% ledger: berger91:readme

\begin{omfigure}
\centering
\begin{tikzpicture}[scale=0.82, font=\scriptsize, >=Stealth]
\draw[thick, rotate=101.37] (-0.72,0) ellipse (2.4 and 2.2895);
\draw[densely dashed, black!60] (-0.331,1.647) -- (0.615,-3.059);
\draw[->, black!70] (0,0) -- (3.9,0);
\node[anchor=south, align=center] at (3.75,0.06) {$\gamma$, towards the\\vernal equinox};
\draw[->, omDef] (0.75,0) arc[start angle=0, end angle=101.37, radius=0.75];
\node[omDef] at (0.78,0.62) {$\varpi$};
\fill[ocRed] (0,0) circle (0.13);
\node[anchor=north west] at (0.05,-0.08) {Sun};
\foreach \x/\y in {-2.062/0, 0/-3.094, 2.321/0, 0/1.688} {\fill[ocBlue] (\x,\y) circle (0.08);}
\node[anchor=east, align=right] at (-2.15,0) {March\\equinox};
\node[anchor=north west, align=left] at (2.38,-0.08) {September\\equinox};
\node[anchor=west] at (0.1,1.92) {December solstice};
\node[anchor=east] at (-0.12,-3.25) {June solstice};
\fill[black] (-0.331,1.647) circle (0.06);
\node[anchor=east] at (-0.42,1.72) {perihelion};
\fill[black] (0.615,-3.059) circle (0.06);
\node[anchor=west] at (0.7,-3.15) {aphelion};
\draw[->, thick, black!60] (-1.9,-1.75) arc[start angle=220, end angle=245, radius=2.6];
\end{tikzpicture}
\par\smallskip{\small The Earth's orbit seen from the north side, with the eccentricity exaggerated
to 0.3 and today's \omterm{def:b2:orbital-forcing:elements}{longitude of perihelion}, $\varpi = \ang{101.4}$. The Earth moves anticlockwise
(arrow); the points mark its position at the equinoxes and solstices. The perihelion falls a few
days after the December solstice, the aphelion a few days after the June solstice. Schematic.}
\end{omfigure}
% ledger: berger91:rows, nasa:milank-detail

Seen from the Earth, the Sun moves along the ecliptic, and its true longitude $\lambda$ is the
angle from the vernal equinox: $\lambda = \ang{90}$ at the June solstice, \ang{270} at the December
solstice. The Earth is at perihelion when the Sun is in the direction opposite to the perihelion,
at $\lambda = \varpi + \ang{180}$; today $\varpi = \ang{101.37}$ and the perihelion falls at
$\lambda = \ang{281.37}$, about eleven days after the December solstice, on or about 3 January.
% ledger: berger91:rows, nasa:milank-detail

\begin{definition}[Milankovitch parameters]\label{def:b2:orbital-forcing:milankovitch}
The \emph{Milankovitch parameters}\index{Milankovitch parameters} are the three orbital
quantities that set the distribution of sunlight over the Earth through the seasons and the
latitudes: the eccentricity $e$, the obliquity $\varepsilon$ and the \omterm{def:b2:orbital-forcing:elements}{climatic precession}
$e\sin\varpi$. The \emph{Milankovitch theory}\index{Milankovitch theory} states that their slow
variations pace the growth and decay of the Quaternary ice sheets, through the summer insolation
at high northern latitudes.
\end{definition}
% ledger: nasa:milank-detail, eo00:milankovitch

Each parameter varies with its own rhythms. The eccentricity has main periods near \num{95000},
\num{99000}, \num{124000} and \num{131000} years, which combine into a cycle of about
\num{100000} years, and a period of \num{405000} years, the most stable of all; over the last
million years it ranged from \num{0.0046} to \num{0.0565}, and today, at \num{0.0172}, it is slowly
decreasing. The obliquity varies with a main period of \num{40978} years, between \ang{22.1} and
\ang{24.5} over the last million years; at \ang{23.44} it is decreasing, past its last maximum
about \num{10000} years ago. The axis wobbles round in about \num{25770} years, while the
orbit turns in about \num{112000}; together they move the perihelion through the seasons with
main periods near \num{19000}, \num{22400} and \num{23700} years.
% ledger: la04:eccentricity, la04:obliquity, la04:precession, nasa:milank-detail, data:b4-orbital

\begin{omfigure}
\centering
% figdata/bachelor-2/orbital.py part elements
\begin{tikzpicture}
\begin{groupplot}[group style={group size=1 by 3, vertical sep=0.3cm, xlabels at=edge bottom,
  xticklabels at=edge bottom}, width=0.96\linewidth, height=3.1cm, xmin=0, xmax=1000,
  xtick={0,100,...,1000}, tick label style={font=\scriptsize}, label style={font=\small},
  scaled y ticks=false, yticklabel style={/pgf/number format/fixed,
  /pgf/number format/precision=2}]
\nextgroupplot[ylabel={$e$}, ymin=0, ymax=0.085, ytick={0,0.02,0.04,0.06}]
\addplot[omDef, thin] table[x=age_ka, y=ecc] {figdata/out/bachelor-2-orbital-elements.dat};
\node[font=\scriptsize, anchor=north west] at (rel axis cs:0.01,0.97) {a\enspace eccentricity};
\nextgroupplot[ylabel={$\varepsilon$ (\unit{\degree})}, ymin=21.9, ymax=25.4, ytick={22,23,24}]
\addplot[ocBlue, thin] table[x=age_ka, y=obl] {figdata/out/bachelor-2-orbital-elements.dat};
\node[font=\scriptsize, anchor=north west] at (rel axis cs:0.01,0.97) {b\enspace obliquity};
\nextgroupplot[ylabel={$e\sin\varpi$}, ymin=-0.07, ymax=0.12, ytick={-0.05,0,0.05},
  xlabel={age (thousands of years before 1950)}]
\addplot[ocRed, thin] table[x=age_ka, y=prec] {figdata/out/bachelor-2-orbital-elements.dat};
\node[font=\scriptsize, anchor=north west] at (rel axis cs:0.01,0.97) {c\enspace climatic
precession};
\end{groupplot}
\end{tikzpicture}
\par\smallskip{\small The three \omterm{def:b2:orbital-forcing:milankovitch}{Milankovitch parameters} over the last million years: (a) the
eccentricity; (b) the obliquity; (c) the \omterm{def:b2:orbital-forcing:elements}{climatic precession}, whose envelope is the eccentricity.
Data: Berger and Loutre (1991), NOAA NCEI Paleoclimatology, as fetched on 23 October 2026.}
\end{omfigure}
% ledger: data:b4-orbit91

The \omterm{def:b2:orbital-forcing:elements}{climatic precession} swings fast, but only as far as the eccentricity lets it: when the orbit is
nearly circular, as \num{400000} years ago, it hardly matters where the perihelion falls. The
figure shows it plainly: the envelope of panel (c) is panel (a).
% ledger: data:b4-orbit91

\section{Sunlight from the orbit}\label{sec:b2:orbital-forcing:insolation}

The daily-mean insolation recalled above needs one more ingredient: the distance of the Sun on a
given day.

\begin{proposition}[Distance and the Sun's longitude]\label{prop:b2:orbital-forcing:distance}
When the Sun's true longitude is $\lambda$, the Earth--Sun distance is
\[
  \frac{r}{a} = \frac{1 - e^2}{1 - e\cos(\lambda - \varpi)} .
\]
At the June solstice ($\lambda = \ang{90}$), to first order in $e$,
$(a/r)^2 \approx 1 - 2e\sin\varpi$; at the December solstice, $(a/r)^2 \approx 1 + 2e\sin\varpi$.
\end{proposition}

\begin{proof}
On an ellipse with the Sun at a focus, $r/a = (1 - e^2)/(1 + e\cos v)$, where the true anomaly $v$
is the angle at the Sun from the perihelion to the Earth. The Earth is seen from the Sun in the
direction opposite to that of the Sun seen from the Earth, so its longitude seen from the Sun is
$\lambda + \ang{180}$, and $v = \lambda + \ang{180} - \varpi$, $\cos v = -\cos(\lambda - \varpi)$.
At $\lambda = \ang{90}$, $\cos(\lambda - \varpi) = \sin\varpi$ and $(a/r)^2 = (1 -
e\sin\varpi)^2/(1 - e^2)^2 = 1 - 2e\sin\varpi + O(e^2)$; at $\lambda = \ang{270}$ the sign changes.
\end{proof}

The \omterm{def:b2:orbital-forcing:elements}{climatic precession} is therefore the right index: it measures how much nearer to the Sun, or
farther, the Earth is at the June solstice. When $e\sin\varpi$ is large and positive, as today and
much more so \num{116000} years ago, the northern summer falls near aphelion and is weak; when it is
large and negative, as about \num{10000} years ago, the northern summer falls near perihelion.
% ledger: berger91:rows

\begin{example}[The June sun at \ang{65} N, today]\label{ex:b2:orbital-forcing:today}
Today $e = 0.017236$, $\varpi = \ang{101.37}$, $\varepsilon = \ang{23.446}$, and the tables use
$S_0 = \qty{1360}{W.m^{-2}}$. At the June solstice $\delta = \varepsilon$; at $\varphi =
\ang{65}$, $\cos h_0 = -\tan\ang{65}\tan\ang{23.446} = -0.9301$, so $h_0 = \ang{158.44} =
\qty{2.7654}{rad}$: the Sun is up for $2 \times 158.44/360 \times 24 = 21.1$~hours. The bracket of
the insolation formula is $2.7654 \times 0.9063 \times 0.3979 + 0.4226 \times 0.9174 \times 0.3674 =
1.1397$, and at the mean distance $\bar Q = 432.90 \times 1.1397 = \qty{493.4}{W.m^{-2}}$. The
distance factor is $(a/r)^2 = 0.9671$, so $\bar Q = \qty{477.1}{W.m^{-2}}$.
\end{example}
% ledger: berger91:rows, berger91:readme, data:b4-orbital

Kepler's second law adds a twist. Near aphelion the Earth moves more slowly: the season that falls
there is weaker but lasts longer. Northern summers are today about 4.5 days longer than northern
winters, and springs about three days longer than autumns.
% ledger: nasa:milank-detail

\begin{proposition}[Annual means do not depend on precession]\label{prop:b2:orbital-forcing:annual}
The energy received per unit area at the top of the atmosphere between two given longitudes of
the Sun does not depend on $\varpi$, and the annual-mean insolation at any latitude is that of a
circular orbit of radius $a$ multiplied by $1/\sqrt{1 - e^2}$. At a pole the annual mean is
\[
  \bar Q_{\text{pole}} = \frac{S_0\sin\varepsilon}{\pi\sqrt{1 - e^2}} ,
\]
which grows with the obliquity.
\end{proposition}

\begin{proof}
By Kepler's second law the radius sweeps area at the constant rate $\pi ab/T$, and the area swept
while the Earth turns by $\mathrm d\lambda$ is $\frac12 r^2\,\mathrm d\lambda$, so $\mathrm dt =
T r^2\,\mathrm d\lambda/(2\pi ab)$, with $b = a\sqrt{1 - e^2}$. The insolation at a given latitude
and longitude of the Sun is $(a/r)^2$ times its value $Q_a(\lambda)$ at the distance $a$, so the
energy received is $\int Q_a(\lambda)(a/r)^2\,\mathrm dt = \frac{T}{2\pi\sqrt{1 - e^2}}\int
Q_a(\lambda)\,\mathrm d\lambda$: the distance has cancelled, and with it $\varpi$. Over the whole
year, $\lambda$ runs through \ang{360} and the mean is $\frac{1}{2\pi\sqrt{1-e^2}}\int_0^{2\pi}
Q_a\,\mathrm d\lambda$. At the North Pole $Q_a = S_0\sin\delta = S_0\sin\varepsilon\sin\lambda$
for $0 < \lambda < \pi$ and 0 otherwise; the integral is $2S_0\sin\varepsilon$.
\end{proof}
% ledger: earth:TSI

Precession moves sunlight from one season to another, never in or out of a year; eccentricity
changes the yearly total by a factor $1/\sqrt{1 - e^2}$, at most \qty{0.16}{\percent} over the
last million years; obliquity moves sunlight between latitudes, raising the annual mean at the
poles and lowering it slightly at the equator. Only the seasonal and latitudinal pattern changes
much, and that is what the figure below shows.
% ledger: data:b4-orbital, nasa:milank-detail

\begin{omfigure}
\centering
% figdata/bachelor-2/orbital.py part lat
\begin{tikzpicture}
\begin{groupplot}[group style={group size=2 by 1, horizontal sep=1.6cm},
  tick label style={font=\scriptsize}, label style={font=\small}, height=5.4cm,
  width=0.48\linewidth, xmin=-90, xmax=90, xtick={-90,-60,...,90}, xlabel={latitude
  (\unit{\degree})}, ylabel={$\bar Q$ (\unit{W.m^{-2}})}]
\nextgroupplot[ymin=0, ymax=650, title={a\enspace June solstice}, title style={font=\scriptsize},
  legend style={font=\scriptsize, at={(0.03,0.97)}, anchor=north west, draw=none},
  legend cell align=left]
\addplot[black, thick] table[x=lat, y=jun_0] {figdata/out/bachelor-2-orbital-lat.dat};
\addlegendentry{today}
\addplot[ocRed, thick, densely dashed] table[x=lat, y=jun_10] {figdata/out/bachelor-2-orbital-lat.dat};
\addlegendentry{\num{10000} years ago}
\nextgroupplot[ymin=100, ymax=500, title={b\enspace annual mean}, title style={font=\scriptsize},
  legend style={font=\scriptsize, at={(0.5,0.03)}, anchor=south, draw=none},
  legend cell align=left]
\addplot[ocBlue, thick] table[x=lat, y=ann_lo] {figdata/out/bachelor-2-orbital-lat.dat};
\addlegendentry{$\varepsilon = \ang{22.1}$}
\addplot[ocRed, thick, densely dashed] table[x=lat, y=ann_hi] {figdata/out/bachelor-2-orbital-lat.dat};
\addlegendentry{$\varepsilon = \ang{24.5}$}
\end{groupplot}
\end{tikzpicture}
\par\smallskip{\small (a) Daily-mean insolation at the June solstice, today and \num{10000} years
ago, when the northern summer fell near perihelion and the obliquity was larger: about
\qty{50}{W.m^{-2}} more in the northern high latitudes. (b) Annual-mean insolation for the two
extreme obliquities of the last million years, with today's eccentricity: a larger tilt sends
more sunlight to the poles and slightly less to the tropics. Computed from the orbital elements
with $S_0 = \qty{1360}{W.m^{-2}}$.}
\end{omfigure}
% ledger: data:b4-orbital, data:b4-orbit91

\section{Computing the insolation curve}\label{sec:b2:orbital-forcing:curve}

The insolation curve that Milankovitch built is now a routine calculation. The orbital elements of
the past are computed by integrating the motions of the planets; tabulated every thousand years,
they are the input. For each age, the method below gives the insolation at any latitude and season.

\begin{method}[Insolation of the past]\label{met:b2:orbital-forcing:curve}
\begin{enumerate}
  \item Read $e$, $\varpi$ and $\varepsilon$ for the age wanted.
  \item Choose the season by the Sun's true longitude $\lambda$: \ang{90} for the June solstice;
        the tables' mid-July corresponds to $\lambda = \ang{120}$.
  \item Compute $\delta$ from $\sin\delta = \sin\varepsilon\sin\lambda$, and the hour angle of
        sunset $h_0$ at the latitude $\varphi$ (polar day or night if $|\tan\varphi\tan\delta| > 1$).
  \item Compute the distance factor $(a/r)^2$ with \cref{prop:b2:orbital-forcing:distance}.
  \item Multiply: $\bar Q = (S_0/\pi)(a/r)^2(h_0\sin\varphi\sin\delta +
        \cos\varphi\cos\delta\sin h_0)$.
\end{enumerate}
\end{method}
% ledger: berger91:readme, data:b4-orbital

Applied to the 1001 ages of the last million years, the method reproduces the tabulated mid-July
insolation at \ang{65} N to within \qty{0.012}{W.m^{-2}}, the rounding of the tables, and the
tabulated June values at the present day exactly. The tables put the origin of time at 1950.
% ledger: data:b4-orbital, berger91:readme

\begin{omfigure}
\centering
% figdata/bachelor-2/orbital.py part jul65
\begin{tikzpicture}
\begin{groupplot}[group style={group size=1 by 2, vertical sep=0.3cm, xlabels at=edge bottom,
  xticklabels at=edge bottom}, width=0.96\linewidth, xmin=0, xmax=400, xtick={0,50,...,400},
  tick label style={font=\scriptsize}, label style={font=\small}]
\nextgroupplot[height=5.0cm, ymin=370, ymax=590, ylabel={$\bar Q$ at \ang{65} N
  (\unit{W.m^{-2}})}, legend style={font=\scriptsize, at={(0.5,1.02)}, anchor=south, draw=none,
  legend columns=3, /tikz/every even column/.append style={column sep=6pt}},
  legend cell align=left]
\addplot[black!45, line width=2.2pt] table[x=age_ka, y=tab] {figdata/out/bachelor-2-orbital-jul65.dat};
\addlegendentry{mid-July, tabulated}
\addplot[ocRed, thin] table[x=age_ka, y=calc] {figdata/out/bachelor-2-orbital-jul65.dat};
\addlegendentry{mid-July, computed}
\addplot[ocBlue, thin, densely dashed] table[x=age_ka, y=jun65] {figdata/out/bachelor-2-orbital-jul65.dat};
\addlegendentry{June solstice, computed}
\nextgroupplot[height=2.8cm, ymin=-0.02, ymax=0.02, ytick={-0.01,0,0.01},
  scaled y ticks=false, yticklabel style={/pgf/number format/fixed}, ylabel={residual},
  xlabel={age (thousands of years before 1950)}]
\addplot[only marks, mark=*, mark size=0.35pt, black] table[x=age_ka, y=resid]
  {figdata/out/bachelor-2-orbital-jul65.dat};
\end{groupplot}
\end{tikzpicture}
\par\smallskip{\small Top: daily-mean insolation at \ang{65} N over the last \num{400000} years,
mid-July as tabulated by Berger and Loutre (grey, wide) and as computed by the method above (red),
and at the June solstice (blue, dashed). Bottom: computed minus tabulated mid-July values, in
\unit{W.m^{-2}}.}
\end{omfigure}
% ledger: data:b4-orbit91, data:b4-orbital

Two features stand out. The curve is dominated by the \omterm{def:b2:orbital-forcing:elements}{climatic precession}, with peaks some
\num{20000} years apart whose height swells and shrinks with the eccentricity; the obliquity adds
a slower modulation. And the swings are large: over the last \num{800000} years the mid-July value
ranged from 388 to \qty{497}{W.m^{-2}}.
% ledger: data:b4-orbital

\section{Milankovitch theory and the glacial cycles}\label{sec:b2:orbital-forcing:theory}

\begin{history}[Sunshine at sixty-five degrees north]
Milutin Milankovitch, born in 1879, built a mathematical theory of climate from the seasonal and latitudinal variations of sunlight. Combining the three orbital variations,
he computed the insolation of the latitudes and the corresponding surface temperatures for
\num{600000} years before 1800. At the suggestion of the climatologist Köppen he chose the summer
insolation at \ang{65} N as the decisive quantity: the great ice sheets had grown near that
latitude, and cool summers melt less of the winter's snow, so that snow accumulates from year to
year. He thought the obliquity the most important of the three parameters. For about fifty years
the theory was largely ignored; Milankovitch died in 1958. In 1976 Hays, Imbrie and Shackleton
found in deep-sea sediments of the last \num{450000} years spectral peaks of climate variability
at about \num{23000}, \num{42000} and \num{100000} years, the periods of precession, obliquity and
eccentricity, and concluded that changes in the Earth's orbit are the fundamental cause of the
Quaternary ice ages.
\end{history}
% ledger: eo00:milankovitch, nasa:milank-detail, hays76:periods

The figure below repeats their test on the records of \cref{ch:b2:climate-proxies}. It compares the
amplitude spectra, the strength of each period, of the \ang{65} N July insolation and of the
benthic stack, in two windows of \num{700000} years.

\begin{omfigure}
\centering
% figdata/bachelor-2/orbital.py part spectrum
\begin{tikzpicture}
\begin{axis}[width=0.9\linewidth, height=5.6cm, xmin=10, xmax=150, ymin=0, ymax=1.08,
  xlabel={period (thousands of years)}, ylabel={relative amplitude},
  xtick={19,23,41,100,150}, tick label style={font=\scriptsize}, label style={font=\small},
  legend style={font=\scriptsize, at={(0.5,1.02)}, anchor=south, draw=none, legend columns=3,
  /tikz/every even column/.append style={column sep=6pt}}, legend cell align=left]
\addplot[black!50, thick] table[x=period, y=amp_ins] {figdata/out/bachelor-2-orbital-spectrum.dat};
\addlegendentry{insolation \ang{65} N, 0--\qty{700}{ka}}
\addplot[ocRed, thick] table[x=period, y=amp_young] {figdata/out/bachelor-2-orbital-spectrum.dat};
\addlegendentry{LR04, 0--\qty{700}{ka}}
\addplot[ocBlue, thick, densely dashed] table[x=period, y=amp_old] {figdata/out/bachelor-2-orbital-spectrum.dat};
\addlegendentry{LR04, \numrange{1250}{1950}~ka}
\end{axis}
\end{tikzpicture}
\par\smallskip{\small Amplitude spectra, each scaled to its own maximum, of the mid-July insolation
at \ang{65} N over the last \num{700000} years and of the LR04 benthic stack over the same window
and over \num{1250000} to \num{1950000} years ago (trend removed, tapered). The insolation has its
power at the precession periods, near 19 to 24 thousand years, and at the obliquity period; the
recent ice volume peaks near \num{100000} years, where the insolation has almost nothing; the older
ice volume peaks at \num{41000} years.}
\end{omfigure}
% ledger: data:b4-orbital, data:lr04, hays76:periods

\subsection*{Glacials, interglacials and terminations}

\begin{definition}[Glacial period; interglacial; glacial termination]\label{def:b2:orbital-forcing:glacials}
In the Quaternary record, a \emph{glacial period}\index{glacial period} is an interval when large
ice sheets cover the northern continents, recorded by high benthic $\delta^{18}$O and low sea level;
an \emph{interglacial}\index{interglacial} is an interval with little land ice in the Northern
Hemisphere outside Greenland, like the present. A \emph{glacial
termination}\index{glacial termination} is one of the major transitions that end a glacial cycle,
from full glacial to interglacial conditions, much faster than the growth that preceded it.
\end{definition}
% ledger: kohler20:mpt, barker21:terminations, ipcc:ref-periods

Over the last half million years the dominant period of climate is close to \num{100000} years,
although the forcing at that period, which comes from the eccentricity alone, is negligible
beside that of obliquity and precession. The cycles are also lopsided: the benthic $\delta^{18}$O
records a stepwise build-up of ice over about \num{90000} years, then a decay in about
\num{10000}, a saw-tooth whose falling edges are the terminations. The last \omterm{def:b2:orbital-forcing:glacials}{interglacial}, from
about \num{129000} to \num{116000} years ago, ended as the summer insolation at \ang{65} N fell
to one of its lowest values of the last \num{800000} years: the June-solstice value reached its
minimum, \qty{438.6}{W.m^{-2}}, \num{116000} years ago, and the mid-July value its second-lowest
minimum, \qty{392.2}{W.m^{-2}}, \num{114000} years ago. The Last Glacial Maximum came some
\num{95000} years later.
% ledger: barker21:terminations, ipcc:ref-periods, data:b4-orbital, data:b4-orbital-extremes

How quickly can an ice sheet follow its forcing? The simplest model treats the ice volume $V$ as
relaxing towards the volume $V_{\mathrm{eq}}$ that the present insolation would sustain, with a
response time $\tau$:
\[
  \frac{\mathrm dV}{\mathrm dt} = \frac{V_{\mathrm{eq}}(t) - V}{\tau} .
\]

\begin{proposition}[A lagged response]\label{prop:b2:orbital-forcing:lag}
For a forcing $V_{\mathrm{eq}} = A\sin\omega t$ the solution of the lag model tends to
\[
  V = \frac{A}{\sqrt{1 + \omega^2\tau^2}}\,\sin\bigl(\omega(t - t_{\ell})\bigr),\qquad
  t_{\ell} = \frac{\arctan(\omega\tau)}{\omega} :
\]
the ice follows slow forcing almost fully and fast forcing only weakly, always late. The response
to a sum of sinusoids is the sum of the responses: the model creates no period that its forcing
lacks.
\end{proposition}

\begin{proof}
Try $V = B\sin\omega t + C\cos\omega t$: substituting, $\omega B\tau = -C$ and $-\omega C\tau = A -
B$, so $B = A/(1 + \omega^2\tau^2)$, $C = -\omega\tau A/(1 + \omega^2\tau^2)$, an amplitude
$\sqrt{B^2 + C^2} = A/\sqrt{1 + \omega^2\tau^2}$ and a phase delay $\arctan(\omega\tau)$. Any other
solution differs from this one by a multiple of $\mathrm e^{-t/\tau}$, which dies away. The
equation is linear, so solutions add.
\end{proof}

\begin{omfigure}
\centering
% figdata/bachelor-2/orbital.py part lag
\begin{tikzpicture}
\begin{axis}[width=0.96\linewidth, height=5.4cm, xmin=0, xmax=400, ymin=-3.4, ymax=4.2,
  xtick={0,50,...,400}, xlabel={age (thousands of years before 1950)},
  ylabel={standardised}, tick label style={font=\scriptsize}, label style={font=\small},
  legend style={font=\scriptsize, at={(0.5,1.02)}, anchor=south, draw=none, legend columns=3,
  /tikz/every even column/.append style={column sep=6pt}}, legend cell align=left]
\addplot[black!35, thin] table[x=age_ka, y=forcing] {figdata/out/bachelor-2-orbital-lag.dat};
\addlegendentry{minus insolation \ang{65} N}
\addplot[ocRed, thick] table[x=age_ka, y=model] {figdata/out/bachelor-2-orbital-lag.dat};
\addlegendentry{lag model, $\tau = \qty{10}{kyr}$}
\addplot[black, thick] table[x=age_ka, y=lr04] {figdata/out/bachelor-2-orbital-lag.dat};
\addlegendentry{LR04 (ice up)}
\end{axis}
\end{tikzpicture}
\par\smallskip{\small The lag model driven by the mid-July insolation at \ang{65} N (less
sunshine, more ice), against the benthic stack, all three standardised (mean 0, standard deviation
1); positive values mean more ice. The model follows the precession wiggles of the stack but has
none of its large saw-tooth cycles. The response time of \qty{10}{kyr} is a model choice.}
\end{omfigure}
% ledger: data:b4-orbital, data:lr04

Over the last \num{800000} years, the lagged model correlates with the stack at $r = 0.42$, the
insolation itself at only $0.23$: the lag helps, but a linear ice sheet cannot make
\num{100000}-year cycles out of a forcing that has none. The terminations must involve
non-linear behaviour that amplifies a modest push from the insolation. The candidates are feedbacks that act on the whole climate:
the \omterm{def:b2:precambrian-earth:ice-albedo}{ice--albedo feedback} of \cref{def:b2:precambrian-earth:ice-albedo}, since a planet with more
snow and ice reflects more sunlight; and carbon dioxide, which in the Antarctic ice cores rises on
time scales of millennia whenever the overturning circulation of the Atlantic is weakened, as it
is when meltwater from the ice sheets pours into the North Atlantic. Beyond a threshold of carbon
dioxide, \omterm{def:b2:orbital-forcing:glacials}{interglacial} conditions can stabilise.
% ledger: data:b4-orbital, barker21:terminations, eo00:milankovitch

\begin{definition}[Mid-Pleistocene transition]\label{def:b2:orbital-forcing:mpt}
The \emph{mid-Pleistocene transition}\index{mid-Pleistocene transition} is the change, from about
\num{1250000} to \num{700000} years ago, from glacial cycles of about \num{41000} years and small
amplitude to quasi-periodic cycles of about \num{100000} years and large amplitude, with thicker
ice sheets, without any corresponding change in the orbital forcing.
\end{definition}
% ledger: clark06:mpt, nasa:milank-detail, kohler20:mpt

Before the transition the ice sheets followed the obliquity, as the dashed spectrum of the figure
above shows; one count finds a \omterm{def:b2:orbital-forcing:glacials}{glacial termination} in 67, 88 and \qty{52}{\percent} of the
obliquity cycles of the early, middle and late Quaternary. Orbital theory alone cannot explain the change; the candidates are the dynamics
of the ice sheets themselves, perhaps after the removal of the soft regolith on which they slid,
and a gradual decline of carbon dioxide during the glacial minima. No explanation is yet
generally agreed.
% ledger: kohler20:mpt, nasa:milank-detail, clark06:mpt

\section{Carbon dioxide in the glacial cycles}\label{sec:b2:orbital-forcing:co2}

The Antarctic ice cores show that carbon dioxide follows the glacial cycles: between 188 and
\qty{194}{ppm} at the Last Glacial Maximum, between 266 and \qty{282}{ppm} in the last \omterm{def:b2:orbital-forcing:glacials}{interglacial};
the composite record of \cref{ch:b2:climate-proxies} gives \qty{189.7}{ppm} for 19 to
\qty{23}{ka} and \qty{279.4}{ppm} over the two thousand years before industry. Its radiative forcing
grows close to the logarithm of the concentration, \qty{3.93}{W.m^{-2}} for each doubling.
% ledger: ipcc:fig2-34, data:b4-proxies-co2, ipcc:log-co2, ipcc:erf-2xco2

\begin{example}[Orbit against carbon dioxide]\label{ex:b2:orbital-forcing:co2}
From \qty{189.7}{ppm} to \qty{279.4}{ppm}, the forcing is $3.93\log_2(279.4/189.7) = 3.93 \times
0.559 = \qty{2.2}{W.m^{-2}}$, over the whole globe and the whole year. The orbit, over the last
million years, changed the global and annual mean of the insolation by at most $340.54 - 340.00 =
\qty{0.54}{W.m^{-2}}$, but the mid-July insolation at \ang{65} N by about \qty{109}{W.m^{-2}}.
\end{example}
% ledger: data:b4-proxies-co2, ipcc:erf-2xco2, ipcc:log-co2, data:b4-orbital

The orbit is a pacemaker, not an engine. It redistributes sunlight between seasons and latitudes
and sets the timing; the ice sheets, by their albedo and their slowness, and carbon dioxide, by
its global forcing, turn a regional signal into a global change of climate.
% ledger: hays76:periods, barker21:terminations, eo00:milankovitch

\section*{Exercises}\label{sec:b2:orbital-forcing:exercises}

Data for the exercises, unless stated otherwise: $S_0 = \qty{1360}{W.m^{-2}}$ (the value of the
tables); a tropical year of \num{365.2422} days; the orbital elements of the tables, of which today
$e = 0.017236$, $\varpi = \ang{101.37}$, $\varepsilon = \ang{23.446}$.
% ledger: berger91:readme, berger91:rows

\begin{exercise}[$\star$]\label{exo:b2:orbital-forcing:1}
Take $e = 0.0167$. (a) Compute the ratio of the aphelion distance $a(1 + e)$ to the perihelion
distance $a(1 - e)$. (b) Express their difference as a percentage of $a$. (c) By what factor is
the sunlight at perihelion stronger than at aphelion?
\end{exercise}
% ledger: earth:ecc, nasa:milank-detail

\begin{exercise}[$\star$]\label{exo:b2:orbital-forcing:2}
The tables give $e = 0.017236$, $\varpi = \ang{101.37}$ today; $e = 0.019971$, $\varpi =
\ang{293.86}$ \num{10000} years ago; $e = 0.019249$, $\varpi = \ang{359.99}$ \num{6000} years ago.
(a) Compute the \omterm{def:b2:orbital-forcing:elements}{climatic precession} for each age. (b) For each, give the Sun's longitude at
perihelion. (c) At which age was the perihelion nearest the June solstice, and at which the
September equinox?
\end{exercise}
% ledger: berger91:rows

\begin{exercise}[$\star$]\label{exo:b2:orbital-forcing:3}
The axis of rotation precesses with a period of about \num{25770} years relative to the stars, while
the orbit turns with a period of about \num{112000} years. Assume that both motions move the
perihelion relative to the equinox in the same sense, so that their frequencies add. (a) Combine
the two frequencies to find the period with which the perihelion
moves through the seasons. (b) Compare with the main periods of the \omterm{def:b2:orbital-forcing:elements}{climatic precession}, \num{19.0},
\num{22.4} and \num{23.7} thousand years.
\end{exercise}
% ledger: nasa:milank-detail, la04:precession

\begin{exercise}[$\star$]\label{exo:b2:orbital-forcing:4}
With \cref{prop:b2:orbital-forcing:annual} and $e = 0.0167$, compute the annual-mean insolation at
the North Pole for obliquities of \ang{22.1} and \ang{24.5}. By what percentage does it grow between
the two?
\end{exercise}
% ledger: nasa:milank-detail, earth:ecc

\begin{exercise}[$\star\star$]\label{exo:b2:orbital-forcing:5}
(a) Compute exactly the distance factor $(a/r)^2$ at the June solstice today and \num{116000} years
ago ($e = 0.043988$, $\varpi = \ang{92.71}$), and compare with the approximation $1 -
2e\sin\varpi$. (b) Compute both at the December solstice. (c) Why is the approximation worse \num{116000}
years ago?
\end{exercise}
% ledger: berger91:rows

\begin{exercise}[$\star\star$]\label{exo:b2:orbital-forcing:6}
Reproduce \cref{ex:b2:orbital-forcing:today} step by step: the hour angle of sunset at \ang{65} N
at the June solstice, the length of the day, the insolation at the mean distance and the
insolation of today.
\end{exercise}
% ledger: berger91:rows

\begin{exercise}[$\star\star$]\label{exo:b2:orbital-forcing:7}
(a) Using $\mathrm dt = T(r/a)^2\,\mathrm d\lambda/(2\pi\sqrt{1 - e^2})$ and $(r/a)^2 \approx 1 +
2e\cos(\lambda - \varpi)$, show that the half-year from the March to the September equinox lasts
about $T/2 + (2T/\pi)\,e\sin\varpi$. (b) Compute by how many days it exceeds the other half-year
today. (c) Compare with the 4.5 days by which northern summers exceed winters and the three days by
which springs exceed autumns.
\end{exercise}
% ledger: nasa:milank-detail, berger91:rows

\begin{exercise}[$\star\star$]\label{exo:b2:orbital-forcing:8}
Over the last million years the eccentricity ranged from \num{0.0046} to \num{0.0565}. (a) Compute
the global and annual mean insolation $S_0/(4\sqrt{1 - e^2})$ at both extremes. (b) By how much, in
\unit{W.m^{-2}} and in per cent, did the orbit change it? (c) Compare with the range of the
mid-July insolation at \ang{65} N, 388 to \qty{497}{W.m^{-2}}.
\end{exercise}
% ledger: data:b4-orbital

\begin{exercise}[$\star\star$]\label{exo:b2:orbital-forcing:9}
With the lag model and $\tau = \qty{10}{kyr}$, compute the relative amplitude and the lag of the
ice volume for sinusoidal forcing of periods 23, 41 and \num{100} thousand years. Which band does
the ice sheet favour, and why?
\end{exercise}

\begin{exercise}[$\star\star\star$]\label{exo:b2:orbital-forcing:10}
The mid-July insolation at \ang{65} N has almost no power at \num{100000} years, yet the ice volume
of the last \num{700000} years has most of its power there. (a) Show that if the forcing is $A_1\sin
\omega_1 t + A_2\sin\omega_2 t$, the lag model responds only at $\omega_1$ and $\omega_2$. (b) Show
that a response containing the term $V^2$, for example $\mathrm dV/\mathrm dt = (V_{\mathrm{eq}} -
V)/\tau + kV^2$, can create the frequencies $\omega_1 - \omega_2$ and $\omega_1 + \omega_2$, by
substituting the linear solution into $V^2$. (c) Two precession terms of periods \num{23.7} and
\num{19.0} thousand years beat with what period? What do you conclude about the \num{100000}-year
cycles?
\end{exercise}
% ledger: la04:precession, barker21:terminations

\begin{exercise}[$\star\star\star$]\label{exo:b2:orbital-forcing:11}
At the North Pole, the energy received per square metre during the polar day is $E = \int S_0(a/r)^2
\sin\delta\,\mathrm dt$ over the half-year when $\delta > 0$. (a) Show with the proof of
\cref{prop:b2:orbital-forcing:annual} that $E = S_0T\sin\varepsilon/(\pi\sqrt{1 - e^2})$. (b) Compute
$E$ in \unit{GJ.m^{-2}} for $\varepsilon = \ang{22.1}$ and \ang{24.5}, with $e = 0.0167$ and $T$ =
365.25 days. (c) The precession does not appear in $E$. Why can it nevertheless change the
melting of an ice sheet?
\end{exercise}
% ledger: nasa:milank-detail

\begin{exercise}[$\star\star\star$]\label{exo:b2:orbital-forcing:12}
(a) Compute the radiative forcing of the rise of carbon dioxide from \qty{189.7}{ppm} to
\qty{279.4}{ppm}, with \qty{3.93}{W.m^{-2}} per doubling and a logarithmic law. (b) The ranges
given for the Last Glacial Maximum and the last \omterm{def:b2:orbital-forcing:glacials}{interglacial} are 188 to \qty{194}{ppm} and 266 to
\qty{282}{ppm}: compute the smallest and the largest forcing they allow. (c) Compare with the change
of the global annual insolation of \cref{exo:b2:orbital-forcing:8} and with the swing of the
\ang{65} N summer insolation. Which is the trigger and which the amplifier?
\end{exercise}
% ledger: data:b4-proxies-co2, ipcc:fig2-34, ipcc:erf-2xco2, ipcc:log-co2, data:b4-orbital

\begin{problem}[{Weekend problem --- Sunlight at \ang{65} N}]\label{pb:b2:orbital-forcing:1}
The last \omterm{def:b2:orbital-forcing:glacials}{interglacial} ended about \num{116000} years ago. How much sunshine did the summer of the
latitude of the great northern ice sheets receive then, compared with today? Use the data of the
exercises. The tables give, \num{116000} years ago, $e = 0.0440$, $\varpi = \ang{92.7}$,
$\varepsilon = \ang{22.52}$; today, rounded in the same way, $e = 0.0172$, $\varpi = \ang{101.4}$,
$\varepsilon = \ang{23.45}$.

\medskip\noindent\textbf{Part I --- The orbit then and now.}
\begin{enumerate}
  \item Compute the \omterm{def:b2:orbital-forcing:elements}{climatic precession} then and now.
  \item Compute the perihelion and aphelion distances, in units of $a$, then and now.
  \item At what longitude of the Sun was the Earth at perihelion then, and now? Near which solstice?
  \item By what factor was the sunlight at perihelion stronger than at aphelion, then and now?
  \item Was the obliquity larger or smaller than today? What does that mean for the summers of the
        high latitudes?
  \item Compute the global and annual mean insolation then and now. Is the difference significant?
\end{enumerate}

\medskip\noindent\textbf{Part II --- The distance at the June solstice.}
\begin{enumerate}[resume]
  \item Show from \cref{prop:b2:orbital-forcing:distance} that at the June solstice $r/a = (1 -
        e^2)/(1 - e\sin\varpi)$.
  \item Compute $r/a$ at the June solstice then and now.
  \item Compute $(a/r)^2$ then and now.
  \item Compare with the approximation $1 - 2e\sin\varpi$.
  \item By what percentage was the solstice sunlight weaker than at the mean distance, then and now?
\end{enumerate}

\medskip\noindent\textbf{Part III --- Sunshine at \ang{65} N.}
\begin{enumerate}[resume]
  \item Give the declination of the Sun at the June solstice then and now.
  \item Compute the hour angle of sunset at \ang{65} N on that day, then and now.
  \item Deduce the length of the day.
  \item Compute the daily-mean insolation at the mean distance.
  \item Compute the daily-mean insolation, then and now.
  \item The tables give, for mid-July, 397.60 and \qty{426.76}{W.m^{-2}}. Why are the solstice
        values higher?
  \item Which counted more in the difference between then and now, the obliquity or the distance?
\end{enumerate}

\medskip\noindent\textbf{Part IV --- The ice sheet's answer.}
\begin{enumerate}[resume]
  \item In the lag model with $\tau = \qty{10}{kyr}$, what fraction of a sudden change of
        $V_{\mathrm{eq}}$ has the ice sheet achieved after 5, 10 and \qty{20}{kyr}?
  \item The June insolation at \ang{65} N fell from a maximum about \num{128000} years ago to the
        minimum of \num{116000} years ago, half a precession cycle. For a sinusoid of period
        \qty{23}{kyr}, what fraction of the forcing does the ice sheet follow, and how late?
  \item Why can an ice sheet not reach, within one precession cycle, the size that a low summer
        insolation would sustain?
  \item The ice took some \num{90000} years to build up to the Last Glacial Maximum and about
        \num{10000} years to melt. Why can the lag model not produce that asymmetry?
  \item Carbon dioxide stood between 266 and \qty{282}{ppm} in the last \omterm{def:b2:orbital-forcing:glacials}{interglacial} and between
        188 and \qty{194}{ppm} at the Last Glacial Maximum. How does it help a regional change of
        sunshine become a global change of climate?
  \item To sum up: what was the daily-mean insolation at \ang{65} N on the June solstice when the
        last \omterm{def:b2:orbital-forcing:glacials}{interglacial} ended, \num{116000} years ago?
\end{enumerate}
\end{problem}
% ledger: berger91:rows, berger91:readme, data:b4-orbital, ipcc:ref-periods, ipcc:fig2-34, barker21:terminations
